
\documentclass[a4paper, 11pt]{article}
\usepackage[french]{babel} \frenchbsetup{og=«,fg=»}
\usepackage{fullpage}
\usepackage{graphicx}
\usepackage[T1]{fontenc}
\usepackage[utf8]{inputenc}
\usepackage{xcolor}
\usepackage{amsthm}
\usepackage{amsmath}
\usepackage{amssymb}
\usepackage{stmaryrd}
\usepackage{proof}
\usepackage{booktabs}
\usepackage{listings}
\usepackage{langs}

\definecolor{link_color}{RGB}{0,0,200}

\usepackage[bookmarks=true,colorlinks=true,
    citecolor={link_color}, filecolor={link_color},
    urlcolor={link_color}, linkcolor={link_color},
    pdfstartview={XYZ null null 1.22}, hyperindex=true]{hyperref}



\theoremstyle{definition}
\newtheorem{definition}{Définition}[section]
\newtheorem{theorem}{Théorème}[section]
\newtheorem{lemma}[theorem]{Lemme}
\newtheorem{proposition}[theorem]{Proposition}
\newtheorem{corollary}[theorem]{Corollaire}

\newcommand{\etat}[1]{\langle #1 \rangle}
\newcommand{\etatz}[1]{\langle #1 \rangle_0}
\newcommand{\etatu}[1]{\langle #1 \rangle_1}
\newcommand{\etatd}[1]{\langle #1 \rangle_2}
\newcommand{\etatt}[1]{\langle #1 \rangle_3}
\newcommand{\mcc}{\mathcal{C}}
\newcommand{\mcd}{\mathcal{D}}
\newcommand{\mcl}{\mathcal{L}}
\newcommand{\mce}{\mathcal{E}}
\newcommand{\llb}{\llbracket}
\newcommand{\rrb}{\rrbracket}
\newcommand{\ra}{\rightarrow}
\newcommand{\rsa}{\rightsquigarrow}

\title{Deux modèles, deux mesures}
\author{Ewen DUFOUR\\[1ex]
{\small Encadré par}\\ {\small Thibaut BALABONSKI}}
\date{mars - juillet 2025}
\begin{document}
\maketitle
\begin{abstract}
\end{abstract}
\section{Introduction}
\paragraph{Thèse d'invariance des modèles.}
Énoncée par Peter van Emde Boas et Cees Slot~\cite{slot1988problem}, 
 la thèse d'invariance des modèles affirme qu'il existe une classe de modèles de calcul, 
 et que tous les modèles au sein de cette classe peuvent se simuler entre eux avec un surcoût polynomial en temps et linéaire en espace. 
 En particulier, cette classe contiendrait les machines de Turing, les RAM, etc.
\paragraph{Le $\lambda$-calcul et sa complexité.}
 Le $\lambda$-calcul est un modèle de calcul servant de base aux langages de programmation fonctionnels. 
 Il est donc naturel de nous assurer que ce modèle est bien dans cette classe.
 Bien qu'étant naturelle sur les machines de Turing, 
 les notions de temps et d'espace sur les $\lambda$-termes et leurs réductions ne sont pas immédiates.
\paragraph{Le linear substitution calculus.}
Le Linear Substitution Calculus est une variante du $\lambda$-calcul introduite par Accattoli et co-auteurs~\cite{accattoli2014nonstandard}.
Les termes sont donnés par la grammaire suivante~:
\begin{align*}
    t ::=~x~|~\lambda x.t~|~t~t~|~t[x/t]
\end{align*}
En plus des termes habituels du $\lambda$-calcul, nous ajoutons une construction $t[x/t']$ qui explicite la substitution de $x$ par $t'$ dans $t$.
Nous ajoutons à ces termes les règles de réduction suivantes~:
\begin{align*}
    \mcl[\lambda x.t]~u~&\ra_{d\beta}~\mcl[t[x/u]]\\
    \mcc\llb x\rrb[x/u]~&\ra_{ls}~\mcc\llb u\rrb[x/u]\\
    t[x/u]~&\ra_{gc}~t & x\notin fv(t)
\end{align*}
où $\mcc$ et $\mcl$ sont des contextes construits par les grammaires~:
\[\mcl ::= \square~|~\mcl[x/t]\]
\[\mcc ::= \square~|~\mcc[x/t]~|~t[x/\mcc]~|~\mcc~t~|~t~\mcc~|~\lambda x.\mcc\]  
\paragraph{Objectifs.} 
Nous souhaitons montrer que le LSC est un modèle de calcul raisonnable pour l'espace. 
Cependant, plusieurs difficultés se présentent. 
D'une part, il nous faut définir une mesure d'espace pour le LSC et fixer une stratégie de réduction. 
D'autre part, il faut montrer que le LSC muni de la stratégie et de la mesure peut être simulé par les machines de Turing et inversement avec un surcoût d'espace linéaire. 
% Une autre difficulté se présentant est la possibilité de capturer les classes de complexité sous-linéaires, en particulier LOG-SPACE.
% Afin de pouvoir caractériser cette classe de complexité, les machines de Turing distinguent un ruban d'entrée qu'elles ne comptent pas comme de l'espace utilisé et ne comptent que l'espace de travail. 
%Ce procédé est beaucoup moins naturel sur les $\lambda$-termes, car ils représentent à la fois le~\og programme~\fg, les~\og  données d'entrée~\fg~et les~\og calculs~\fg~en cours. 
\paragraph{Travaux connexes.}
  Dans leurs travaux, \cite{forster2019weak}, Forster et al. montrent que le $\lambda$-calcul en appel par valeur est raisonnable pour l'espace (et le temps).
  Cependant, leur mesure ne permet pas de capturer les classes sous-linéaires.
  Accattoli, Dal Lago et Vanoni décrivent une manière de capturer ces classes dans leurs travaux~\cite{accattoli2022reasonable}. 
  Néanmoins, leurs mesures se basent sur une machine abstraite et non pas sur les $\lambda$-termes.

\section{Formalismes}
\subsection{LS(G)C}
\subsubsection{Stratégie de réduction}
Pour cette stratégie de réduction, nous distinguons une classe de termes que nous appellerons des valeurs.
Ceux-ci sont donnés par la grammaire~:
\[v ::= \lambda x.t~|~v[x/t]\]
Les valeurs sont donc les abstractions munies d'une \og liste\fg~de substitutions. Nous distinguons clairement ici le concept de fermeture.
De plus, nous restreignons les contextes que nous utiliserons dans nos réductions à ceux donnés par la grammaire : 
\[\mcc ::= \square~|~\mcc[x/t]~|~\mcc~t~|~v~\mcc\]
Nous remarquons que l'on ne permet plus la substitution sous les abstractions ni dans les substitutions explicites.
Notre stratégie est donc définie ainsi~:
\begin{figure}[!h]
    \centering
    \begin{align*}
    &\infer[LA]{t~u\rightarrow t'~u}{t \rightarrow t'}
 &\infer[RA]{v~t\rightarrow v~t'}{t \rightarrow t'}\\
&\infer[GC]{t[x/u]\rightarrow t'}{t \rightarrow t' & x \notin fv(t') }
&\infer[US]{t[x/u]\rightarrow t'[x/u]}{t \rightarrow t' & x \in fv(t')}\\
&\infer[LS]{\mathcal{C}\llbracket x \rrbracket [x/u]\rightarrow \mathcal{C}\llbracket u \rrbracket [x/u]}{x \in fv(\mathcal{C}\llbracket u \rrbracket) }
&\infer[LSGC]{\mathcal{C}\llbracket x \rrbracket [x/u]\rightarrow \mathcal{C}\llbracket u \rrbracket}{x \notin fv(\mathcal{C}\llbracket u \rrbracket) }\\
&\infer[DB]{\mathcal{L}[\lambda x.t]~v \rightarrow \mathcal{L}[t[x/v]]}{x\in fv(t)}
&\infer[DBGC]{\mathcal{L}[\lambda x.t]~v \rightarrow \mathcal{L}[t]}{x\notin fv(t)}
\end{align*}
    \caption{Règles de réduction du LS(G)C}
\end{figure}
\paragraph{Explications.}
Les règles $RA$ et $LA$ sont la base de l'appel par valeur. Nous commençons par évaluer le terme de gauche d'une application jusqu'à obtenir une abstraction avec ou sans environnement, puis on évalue le terme de droite jusqu'à obtenir une valeur. Ensuite, les règles $DB$ et $DBGC$ interviennent. Leur rôle est d'effectuer l'application et de ne garder la substitution explicite que si elle est utile ($DBGC$ étant la composition de $\ra_{d\beta}$ et $\ra_{gc}$). Les règles $LS$ et $LSGC$ servent à effectuer une substitution lorsque l'on évalue une variable et à supprimer la substitution explicite si elle devient inutile. Enfin, les règles $GC$ et $US$ vont aller évaluer les termes sous les substitutions explicites et les supprimer si elles deviennent inutiles après l'évaluation du sous terme. Cela permet de propager le \og gc \fg~et de ne garder que les substitutions utiles.

Nous considérons que les termes sont représentés en utilisant des indices de de Bruijn. Cependant, les indices peuvent faire référence soit aux abstractions, soit aux substitutions explicites. Voici deux exemples de termes avec cette notation~:
\begin{align*}
    (\lambda.0~1)[\lambda.0] &= (\lambda x.x~y)[y/\lambda z.z]\\
    \lambda.(0~1)[\lambda.0] &=\lambda x.(y~x)[y/\lambda z.z]
\end{align*}
De plus, les termes au sein des substitutions explicites peuvent elles-mêmes faire référence à d'autres substitutions. Il nous faut donc faire attention à bien gérer les indices lors des $DB$, $DBGC$, $LS$, $LSGC$ et $GC$. 
\subsubsection{Mesure d'espace pris par les termes}
Nous définissons la mesure d'espace ainsi~:
\begin{align*}
    |x| &= |dB(x)|\\
    |\lambda x.t| &= 1 + |t|\\
    |t~u| &= |t|+|u|\\
    |t[x/u]| &= |t|+|u|
\end{align*}
L'utilisation d'indices de de Bruijn nous permet de ne pas avoir à stocker les variables des substitutions explicites, ce qui explique leurs coûts. 
Nous étendons naturellement cette définition aux contextes.
\subsection{La machine FLAME}

FLAME est une machine abstraite qui implante la stratégie d'appel par valeur pour les termes du LSC. 
Un état de FLAME se décrit par un triplet annoté de la forme $\etat{(v,n),t,\mcc}_e$ où~:
\begin{enumerate}
  \item $t$ représente le terme en train d'être évaluer.
  \item $\mcc$ est une pile de contexte d'évaluation.
    \item $(v, n)$ contient optionellement l'argument d'une application et sa \og distance\fg~par rapport à l'abstraction.
    \item $e$ est un nombre qui indique la phase de fonctionnement de FLAME.
\end{enumerate} 
Toutes les transitions de FLAME sont décrites dans la figure~\ref{FLAME} en annexe.
\subsubsection{Fonctionnement}
\begin{definition}[Décomposition d'un contexte]
  \label{def:dec}
La décomposition d'un contexte correspond à la liste des parties \og atomiques \fg~qui le composent.
  \begin{align*}
    d(\square) &= \square \\
    d(\mcc[x/v]) &= d(\mcc)::\square[x/v]\\
    d(\mcc~t) &= d(\mcc)::\square~t\\
    d(v~\mcc) &= d(\mcc)::v~\square
  \end{align*} 
Cette définition s'étend naturellement à tout type de contexte.
\end{definition}

L'évaluation des termes se fait en trois phases~: la décomposition, la réduction et la reconstruction. 
\paragraph{Décomposition.}
Dans le premier temps, la machine décompose les termes jusqu'à trouver une réduction possible (substitution ou $\beta$-réduction). 
La machine sépare les applications en évaluant un des membres et stocke l'autre sur sa pile de contextes. 
De même, elle stocke sur sa pile les substitutions et évalue leurs sous termes gauches (si elle ne fait pas partie d'une valeur). 
Nous voyons immédiatement que cette phase implémente la fonction de décomposition d'un contexte (\ref{def:dec}).
Cependant, elle ne stocke pas les contextes vide sur la pile.
Cette phase d'exploration est réalisée par les transitions dont l'état de départ est 0. 

\paragraph{Réduction.}
Lorsque la machine a trouvé une réduction, elle va effectuer une transition $\rsa_{db}$ et rentrer dans l'état 1 si c'est une $\beta$-réduction, 
ou effectuer une transition $\rsa_{ls}$ et entrer dans l'état 2 si c'est une substitution. 
La machine se sert de l'état 1 pour réaliser une réduction du genre $d\beta$, 
il est utilisé pour calculer le décalage des indices de de Bruijn et mettre la nouvelle substitution sous les potentielles substitutions déja présentes. 
La machine entre ensuite dans l'état 2.

\paragraph{Reconstruction.}
Après avoir effectué une réduction, la machine va reconstruire le terme. 
Lors de cette phase, elle récupère les sous termes de sa pile de contexte et reconstruit le terme. 
Lorsqu'elle trouve une substitutions dans sa pile de contexte, elle vérifie que la substituions est bien liée dans son terme. 
Si ce n'est pas le cas elle supprime cette substitution lors de la transition $\rsa_{gc}$. 
Dans le cas contraire, elle replace simplement la substitution sur son terme courant.

\paragraph{Relance.}
Lorsque la machine a fini de reconstruire un terme, si celui-ci est une valeur, elle entre dans l'état 3 et s'arrête complètement. 
Dans le cas contraire, la machine retourne dans l'état 0 et effectue un nouveau cycle d'évaluation.

\subsubsection{Mesure d'espace}
Afin d'effectuer nos analyses sur l'espace pris par notre machine, nous définissons la mesure suivante~:
\begin{align*}
    |\etat{(v,n), t, \mcc}_e| = |t| + |v| + |n| + |\mcc|
\end{align*}
Dans le cas où l'entrée \textbf{arg} de la machine est \textbf{-}, nous prenons comme taille $0$. 
L'entrée \textbf{état} de la machine ne pouvant être que \textbf{0},\textbf{1},\textbf{2} ou \textbf{3} nous ne comptons pas cette information dans nos calculs.
\section{Simulation de LS(G)C}
\subsection{Propreté de LS(G)C}
Étant donné un terme $t$, nous disons de ce terme qu'il est propre s'il ne possède aucune substitution explicite inutile.
Cela nous assure que nous ne pourrons pas effectuer de réduction $GC$ sur le terme $t$.
\begin{definition}[Propreté]
  Un terme $t$ est dit propre si $cl(t)$ possède une dérivation.
\begin{figure}[h!]
  \begin{align*} 
    &\infer{cl(x)}{}
 &\infer{cl(\lambda x.t)}{cl(t)}\\
 &\infer{cl(t~u)}{cl(t) & cl(u) }
&\infer{cl(t[x/u])}{cl(t)&cl(u)&x\in fv(t)}
  \end{align*}
  \caption{Règles de dérivation de la propreté.}\label{fig:clean}
\end{figure}
\end{definition}

\begin{theorem}[Préservation de la propreté]
  Soient $t_0$ et $t_1$ des termes. 
  \[
    t_0~\text{propre}~\wedge~t_0\rightarrow~t_1~\Rightarrow~t_1~\text{propre}
  \]
\end{theorem}
\begin{proof}
  Par induction sur la dérivation de $t_0 \rightarrow t_1$.
\end{proof}

Ce théorème nous assure donc que nos règles de calculs ne produisent aucun terme avec des substitutions explicites pouvant être nettoyées.

\subsection{Comportement de FLAME}
Afin de montrer que FLAME simule raisonablement LS(G)C, nous devons nous intéresser a ses différentes phases.
\begin{lemma}[Décomposition]
  \label{lemmaplouf}
  Soient $\mathcal{C}$, $\mathcal{D}$ des contextes d'évaluation et $t$ un terme.
\[\langle -,\mathcal{C}[t], \mathcal{D}\rangle_0 \rightsquigarrow^* \langle -,t,\mathcal{C} :: \mathcal{D}\rangle_0\]
en utilisant un espace $|\mcc|+|t|+|\mcd|$
\end{lemma}
\begin{proof}
  Par induction sur la forme de $\mathcal{C}$.
\end{proof}
\begin{lemma}[Plongeon sous les substitutions]
  \label{lemmaplouf2}
  Soient $\mathcal{D}$ un contexte d'évaluation, $\mcl$ un contexte de substitution et $t$, $u$ des termes.
  \[ \etatu{(u,n), \mcl[t], \mcd} \rightsquigarrow_{rdb}^k \etatu{(u,n+k), t, \mcl::\mcd} \]
  où $k$ est la longueur de la décomposition de $\mcl$,
en utilisant un espace $|\mcd|+|t|+|\mcl|+|u| + |n+k|$
\end{lemma}
\begin{proof}
  Par induction sur la forme de $\mathcal{L}$.
\end{proof}
\begin{lemma}[Reconstruction]
  \label{lemmaportugal}
  Soient $\mathcal{C}$ et $\mathcal{D}$ des contextes d'évaluation et $t$ un terme.
  \[\mathcal{C}[t]~\text{propre}~\Rightarrow \langle-,t,\mathcal{C}::\mathcal{D}\rangle_2 \rightsquigarrow^* \langle-,\mathcal{C}[t],\mathcal{D}\rangle_2\]
en utilisant un espace $|\mcc|+|t|+|\mcd|$
\end{lemma}
\begin{proof}
  Par induction sur la forme de $\mathcal{C}$.
\end{proof}
Grâce au lemme de décomposition \ref{lemmaplouf}, 
nous nous assurons que la machine effectue correctement la décomposition des contextes et quantifions l'espace pris par cette opération.
Le lemme de plongeon sous les substitutions \ref{lemmaplouf2} 
nous assure quant à lui que l'argument d'une application est bien ammené devant les substitutions explicites lors d'une \og transition d$\beta$ \fg.
De plus, nous voyons que le décalage à effectuer sur les indices de de Bruijn est bien calculé et que nous donnons l'espace utilisé par cette phase.
Finalement, le lemme de reconstruction \ref{lemmaportugal} nous montre que le terme est correctement reconstruit après son évaluation 
et nous donne également l'espace pris par la machine pour cette opération.

\subsection{Simulation raisonnable de LS(G)C par FLAME}
Grâce aux lemmes que nous venons de décrire, nous pouvons donc montrer que FLAME simule raisonnablement LS(G)C.
\begin{theorem}[Simulation raisonnable d'un pas]
  Soient $t$ et $u$ des termes propres et $\mcd$ un contexte d'évaluation.
  Il existe une constante $k$ strictement positive telle que
  \begin{center}  
    $\text{Si}~t \rightarrow u~\text{avec un espace}~S\text{, alors}$\\  
  $\etat{-,t,\mcd}_0 \rightsquigarrow^* \etat{-,u,\mcd}_2~\text{avec un espace}~k\times S+|\mcd|$
  \end{center}
\end{theorem}
\begin{proof}
  Par induction sur la dérivation de $t\rightarrow u$.
\end{proof}

\begin{corollary}
  Soient $t$ et $u$ des termes. Il existe une constante $k$ strictement positive telle que
  \begin{center}
    Si $t \ra^* u$ avec $u$ une valeur en espace $S$, alors\\
    $\etatz{-,t,\square} \rsa^* \etatt{-,u,\square}$ avec un espace $k\times S$
  \end{center}
\end{corollary}

\section{Simulation de FLAME}
Montrons maintenant que FLAME peut être simulée par des machines de Turing avec un surcoût en espace linéaire.
Pour ce faire, il nous faut trois choses~:
\begin{enumerate}
    \item Un encodage des termes ayant un surcoût d'espace linéaire par rapport à la taille de nos termes.
    \item Une représentation des différentes entrées de la machine.
    \item Des fonctions \textbf{red}, \textbf{lift}, \textbf{isFree} et \textbf{isValue}. Ces fonctions vont respectivement réduire ou augmenter les indices de de Bruijn dans un terme encodé, vérifier qu'une variable est libre dans un terme encodé et vérifier qu'un terme encodé est une valeur.
    De nouveau, ces fonctions doivent se faire avec un coût spatial au plus linéaire.
\end{enumerate}

\paragraph{Encodage.} L'encodage des termes se fait via une fonction récursive définie ainsi~:
\label{enc}
\begin{align*}
    \mce(x) &=x\\
    \mce(\lambda .t) &=\lambda \mce(t) A\\
    \mce(t~u)&=(\mce(t)\text{\textvisiblespace}\mce(u))\\
    \mce(t[v])&=S\mce(t)[\mce(v)]
\end{align*}
Nous partons du principe que les termes avant encodage sont sous la forme indice de de Bruijn et les variables sont écrit en base 2. 
Nous avons ajouté un symbole $A$ pour délimiter la fin d'une abstraction et un symbole $S$ pour définir le début d'une substitution explicite. 
Enfin, nous ajoutons des parenthèses et un symbole d'espace aux applications pour délimiter chaque membre.
Il est direct que notre encodage à une taille linéaire en la taille du terme non encodé.
Par la suite, nous noterons $\bar{t}$ l'encodage du terme $t$.
\paragraph{Choix du modèle de calcul.}
Pour représenter les différentes entrées de la machine, nous allons utiliser une machine à quatre rubans. Pour les arguments, nous écrivons l'encodage du terme suivi par le nombre $n$ (la \og distance à l'abstraction) sur le premier ruban. Les termes sont décrits par leur encodage sur le deuxième ruban. Nous sauvegardons les contextes sur me troisième ruban où nous écrivons les encodages des sous-termes précédés par des symboles qui permettent de différencier les types de contexte (LS, RA et LA dans FLAME). Sur le dernier ruban, nous gardons en mémoire l'état de la machine et écrivons les résultats des fonctions \textbf{isFree} et \textbf{isValue}.
Ce choix de modèle fonctionne, car les machines de Turing à plusieurs rubans peuvent être simulées par des machines à un ruban avec un espace linéaire (Théorème 1 des travaux de van Emde Boas \cite{van1987machine}). 
\subsection{Fonctions internes de FLAME}

Maintenant, voyons comment effectuer les différentes opérations que fait la machine sur notre encodage.
Pour plus de clarté, nous allons étudier des programmes qui calculent ces fonctions.

\paragraph{isValue.} Il suffit de parcourir l'encodage de gauche à droite. Si nous rencontrons sur le symbole $S$, nous continuons. Si le symbole rencontré est $\lambda$ renvoyer $vrai$, sinon renvoyer $faux$. Tout cela se fait en espace constant.
\paragraph{isFree.} Le programme en figure \ref{fig:isFree} permet de décider si une variable est libre dans l'encodage d'un terme.
L'algorithme n'accède à l'encodage qu'en lecture et le parcours une unique fois de gauche à droite. Ainsi, celui-ci n'a besoin de garder en mémoire que le nombre de lieurs rencontrés.
Le programme s'exécute donc avec un espace logarithmique en la longueur de l'encodage.
\paragraph{lift.}

Le programme en figure~\ref{fig:lift} permet de calculer l'augmentation des indices de de Bruijn dans un terme encodé.
Pour ce faire, la fonction garde en mémoire l'encodage en augmentant certains indices et change également le nombre de lieurs rencontrés. 
La sauvegarde de l'encodage se fait en espace linéaire en la longueur de l'encodage. 
La taille du nombre de lieurs est quant à elle logarithmique en la longueur de l'encodage.
La fonction s'exécute ainsi dans un espace linéaire en la longueur de l'encodage.
\paragraph{red.}
Pour des raisons similaires à celles données pour la fonction lift, le programme en figure \ref{fig:red} calcule la réduction des indices de de Bruijn libre dans un encodage avec un espace linéaire.\\\\

Les descriptions ci-dessus permettent donc de justifier l'existence d'une machine de Turing simulant FLAME avec un surcoût d'espace linéaire.
\begin{theorem}
  Il existe une machine de Turing $\mathcal{M}$ et une constante $k$ telles que si $\etatz{-, t, \square} \rsa^* \etatt{-, u, \square}$ avec un espace $S$, alors
  $\mathcal{M}~\bar{t} \ra^*~\bar{u}$ avec un espace $k\times S$. 
\end{theorem}
\section{Implémentation}
Afin de pouvoir tester nos machines, nous avons créé Flash, un langage de script permettant de définir des termes et de les réduire.
Un exemple de programme flash se trouve dans la figure \ref{fig:flash} en annexe.
\subsection{Structure d'un programme}
Un programme Flash est une séquence de scripts séparés par \og$;;$\fg.
Chaque script est composé de deux choses~: des définitions de termes et des expressions à évaluer.
Les déclarations et expressions d'un même script sont séparées par \og$;$\fg. 
\paragraph{Définition des termes.}
Nous pouvons définir deux types de termes différents ($\lambda$-termes et lsc-termes).
Les déclarations des termes se font sous leur forme nommée.
Par la suite les $\lambda$-termes seront gardés sous leur forme nommée.
Au contraire, lsc-termes sont représentés sous leurs formes d'indices de de Bruijn et nous annotons les lieurs afin d'améliorer la lisibilité lors de l'affichage des termes.
Cette différence dans la représentation s'explique par le format que les machines que nous voulions tester utilisent. 
En plus de la possibilité de définir des termes, nous avons ajouté des termes prédéfinis.
Nous pouvons également utiliser directement des nombres qui seront encodés avec un encodage de Scott.

\subsubsection{Expressions}
Chaque expression est une suite d'applications qui peuvent soit prendre la forme $e_1~e_2$ (on applique $e_1$ à $e_2$),
ou bien la forme $e_2 |> e_1$ (on passe en argument $e_2$ à $e_1$).
Dans ces expressions, nous pouvons utiliser les termes définis au début du scripts, des machines abstraites, 
la variable red\_size (qui contient l'espace utilisé par la dernière réduction) ou encore des fonctions prédéfinies.  

\paragraph{Machines abstraites.}
Nous avons implémenté deux machines dans flash.
Premièrement, nous avons FLAME \ref{FLAME} qui prend en entrée et renvoie des lsc termes.
Deuxièmement, nous implémentons la Space-KAM \cite{accattoli2022reasonable}, celle-ci va prendre en paramètre et renvoyer des $\lambda$-termes.
Les deux machines mettent à jour red\_size au cours de la réduction des termes.
Nous pouvons optionellement afficher les termes a chaque étapes de réduction des machines.

\paragraph{Fonctions prédéfinies.}
En plus des machines abstraites, nous avons ajouter des fonctions utilitaires qui permettent, par exemple, 
d'afficher des termes ou du texte, calculer la taille d'un terme, de l'encoder(\ref{enc}) ou encore de déplier un lsc-terme.
\section{Conclusion et perspectives}

\section{Remerciements}
Je tiens à remercier très chaleureusement Thibaut pour avoir encadré ce stage ainsi que pour toute l'aide qu'il m'a apporté. 
Je tiens également à remercier Travis pour son aide durant ce stage. Enfin, je tiens à remercier l'équipe du LMF qui s'est montrée chaleureuse et très accueillante.
\bibliographystyle{abbrv}
\bibliography{main}
\section*{Annexe}
\begin{figure}[h]
\label{FLAME}
\centering
\resizebox{\textwidth}{!}{%
\begin{tabular}{@{}l l l l l l l l l l@{}}
\toprule
\textbf{arg} & \textbf{terme} & \textbf{contexte} & \textbf{état} & \textbf{transition}& \textbf{arg} & \textbf{terme} & \textbf{contexte} & \textbf{état} \\
\midrule

– & $t[ u ] $ & $\mcc$ & 0 & $\rsa_{us}$ & – & $t$ & $\textsf{LS}(u) :: \mcc$ & 0  \\

- & $t~u$ & $\mcc$ & 0 & $\rsa_{lapp}$ & – & $t$ & $\textsf{RA}(u) ::\mcc$ & 0  \\
- & $v~t$ & $\mcc$ & 0 & $\rsa_{rapp}$ & – & $t$ & $\textsf{LA}(v) :: \mcc$ & 0  \\
- & $n$  & $\mcc$ & 0 & $\rsa_{ls}$ & - & $\uparrow^{n+1} t$ & $\mcc$ & 2 & $\textsf{LS}(t)=\textsf{nth\_subst}(\mcc)$ \\
- & $v_1~v_2$  & $\mcc$ & 0 & $\rsa_{db}$ & $(v_2,0)$ & $v_1$ & $\mcc$ & 1  \\

$(v_2, n)$ & $v_1[ t ]$ &$\mcc$ & 1 & $\rsa_{rdb}$ & $(v_2, n+1)$ & $v_1$ & $\textsf{LS}(t) :: \mcc$ & 1\\ 
$(v, n)$ & $\lambda.t $ &$\mcc$ & 1 & $\rsa_{fdb}$ & - & $t$ & $\textsf{LS}(\uparrow^{n}v) :: \mcc$ & 2 & $0\in fv(t)$\\
$(v, n)$ & $\lambda.t $ &$\mcc$ & 1 & $\rsa_{fdbgc}$ & - & $\downarrow t$ & $\mcc$ & 2& $0\notin fv(t)$\\
-&$v$ & $[]$ & 2 & $\rsa_{stop}$ & -& $v$ &$[]$ & 3 \\
-&$t$ & $[]$ & 2 & $\rsa_{dive}$ & -& $t$ &$[]$ & 0 \\
-&$t$ & $\textsf{LS}(u) :: \mcc$ & 2 & $\rsa_{gc}$ & -& $\downarrow t$ &$\mcc$ & 2 & $0\notin \textsf{fv}(t)$ \\
-&$t$ & $\textsf{LS}(u) :: \mcc$ & 2 & $\rsa_{bsubst}$ & -& $t[ u ]$ &$\mcc$ & 2 & $0\in \textsf{fv}(t)$ \\
-&$t$ & $\textsf{RA}(u) :: \mcc$ & 2 & $\rsa_{blapp}$ & -& $t~u$ &$\mcc$ & 2 \\
-&$t$ & $\textsf{LA}(u) :: \mcc$ & 2 & $\rsa_{brapp}$ & -& $u~t$ &$\mcc$ & 2 \\

\bottomrule
\end{tabular}
}
\caption{Transitions de FLAME.}
\end{figure}

\begin{figure}[h]
  \begin{center}    
\begin{lstlisting}
let rec free (n : int) (e : encoding) : bool =
  if n < 0 then false
  else
    match e with
    | [] -> false
    | BeginLambda :: el -> free (n + 1) el
    | EndLambda :: el -> free (n - 1) el
    | BeginSubst :: el -> free (n + 1) el
    | (LeftPar  | RightPar  | RightBrack | Space) :: el ->
        free n el
    | LeftBrack :: el -> free (n - 1) el
    | Num i :: el -> if i = n then true else free n el
\end{lstlisting}
  \end{center}
  \caption{Programme OCaml décidant si une variable est libre.}\label{fig:isFree}
\end{figure}

\begin{figure}[h]
  \begin{center}
\begin{lstlisting}
let lift_encoding shift e =
  let rec aux shift bound_nb acc = function
    | [] -> List.rev acc
    | Num n :: es when n < bound_nb -> aux shift bound_nb (Num n :: acc) es
    | Num n :: es -> aux shift bound_nb (Num (n + shift) :: acc) es
    | BeginLambda :: es -> aux shift (bound_nb + 1) (BeginLambda :: acc) es
    | BeginSubst :: es -> aux shift (bound_nb + 1) (BeginSubst :: acc) es
    | LeftBrack :: es -> aux shift (bound_nb - 1) (LeftBrack :: acc) es
    | RightBrack :: es -> aux shift bound_nb (RightBrack :: acc) es
    | Space :: es -> aux shift bound_nb (Space :: acc) es
    | LeftPar :: es -> aux shift bound_nb (LeftPar :: acc) es
    | RightPar :: es -> aux shift bound_nb (RightPar :: acc) es
    | EndLambda :: es -> aux shift (bound_nb - 1) (EndLambda :: acc) es
  in
  aux shift 0 [] e
\end{lstlisting}
  \end{center}
  \caption{Programme OCaml augmentant les indices dans un encodage.}\label{fig:lift}
\end{figure}

\begin{figure}[h]
  \begin{center}
\begin{lstlisting}
let red_encoding e =
  let rec aux n acc = function
    | [] -> List.rev acc
    | Num n' :: es when n' > n -> aux n (Num (n' - 1) :: acc) es
    | Num n' :: es -> aux n (Num n' :: acc) es
    | BeginLambda :: es -> aux (n + 1) (BeginLambda :: acc) es
    | BeginSubst :: es -> aux (n + 1) (BeginSubst :: acc) es
    | LeftBrack :: es -> aux (n - 1) (LeftBrack :: acc) es
    | RightBrack :: es -> aux n (RightBrack :: acc) es
    | Space :: es -> aux n (Space :: acc) es
    | LeftPar :: es -> aux n (LeftPar :: acc) es
    | RightPar :: es -> aux n (RightPar :: acc) es
    | EndLambda :: es -> aux (n - 1) (EndLambda :: acc) es
  in
  aux 0 [] e
\end{lstlisting}
  \end{center}
  \caption{Programme OCaml réduisant les indices dans un encodage.}\label{fig:red}
\end{figure}

\begin{figure}[h]
  \begin{center}
\begin{lstlisting}
(*script 1*)
lsc term t1 = (fun x -> x) ;
lsc term t2 = (fun x y -> y); 

t1 t2 |> flame|> print;
t1 |> print;
t2 |> print;;



(*script2*)
lsc term one = (fun s z -> s (fun s z -> z)); 
lsc term n = 40;
lsc term add = 
  let succ = (fun n s z -> s n) in
  fix (fun f x y -> y (fun n -> succ (f x n)) x) ;

add one one |> print |> flame |> print |> unfold |> print;
red_size |> print;
print "pause\n\n\n";
print (add one one) |> flame |> size |> print;
"" |> print |> print |> print ;
add one n |> encode |> print;
add one n |> flame |> unfold |> print |> encode |> print;
red_size |> print

\end{lstlisting}
  \end{center}
  \caption{Exemple de script Flash.}\label{fig:flash}
\end{figure}
\end{document}
